\chapter{Results} \label{chap: Results}

The simulations discussed in \autoref{sec: performed simulations} are carried out on the Fritz and Helma computation clusters of the \acrfull{nhr}.
Fritz is a CPU cluster that is designed to run parallel jobs.
It is equipped with two Intel Xeon Platinum 8360Y processors per compute node, with $36$ cores each.
To enhance the computation capabilities and available \acrfull{ram} per job, multiple compute nodes can be used concurrently.
For the more demanding simulations involving higher Reynolds numbers, the Helma cluster was utilized.
Helma is a general-purpose \acrshort{gpu} cluster equipped with 4x Nvidia H100 GPUs per compute node.
Like on Fritz, multiple nodes can be used for one job, to improve execution times and enhance the available \acrshort{ram}.

The representation of the data is chosen consistently throughout this chapter.
The three surfaces are differentiated by the colour family of the lines.
As described in \autoref{sec: performed simulations}, $S1$ exhibits the largest roughness height $k$, $S2$ decreases $k$, while keeping the effective slope $Es$ constant and $S3$ shows the same topography as $S1$ but is scaled down, such that it matches the roughness height of $S2$.
The colour convention for each surface at each Reynolds number is shown in \autoref{tab: line style convention}.
The Reynolds numbers are differentiated by the intensity of the colour.
Darker colours represent higher Reynolds numbers.

\begin{table}
    \centering
    \caption[Colour convention for results chapter]{Colour convention used throughout this chapter. Line colour denotes the surface, line style denotes the Reynolds number.}
    \begin{tabular}{c c c c}
        \toprule
        $Re_b$    & S1                        & S2                        & S3 \\
        \midrule
        $5300$   & \sOneFiveThreeSymbol      & \sTwoFiveThreeSymbol      & \sThreeFiveThreeSymbol     \\
        $11700$  & \sOneElevenSevenSymbol    & \sTwoElevenSevenSymbol    & \sThreeElevenSevenSymbol  \\
        $19000$  & \sOneNineteenSymbol       & \sTwoNineteenSymbol       & \sThreeNineteenSymbol   \\
        \bottomrule
    \end{tabular}
    \label{tab: line style convention}
\end{table}

\fxinfo{wurde schon korrigiret}
\section{Simulations with 5300$Re_b$}\label{sec: results 5300Reb}
In this section, the results of the \acrshort{dns} study with the Reynolds number $Re_b = 5300$ are presented.
A comparison is made between the three surfaces.
Additionally, reference data for the smooth pipe at the same Reynolds number is plotted and taken from~\cite{yaoDirectNumericalSimulations2023}.
\begin{table}
    \centering
    \caption[Computed friction values for $Re_b = 5300$]{Computed friction values for $Re_b = 5300$}
    \begin{tabular}{c c c c c c c c c}
        \toprule
                & $Re_\tau$ & $k^+$     & $C_f\ [\times10^{-2}]$ & $\tau_w\ [\times10^{-3}]$ & $\tau_v\ [\times10^{-3}]$ & $\tau_p\ [\times10^{-3}]$ & $\tau_v/\tau_w$   & $\Delta U^+$\\
        \midrule
        smooth  & $181.0$   & -         & $2.32$   & $1.14$   & $1.14$   & -                     & $1.00$           & - \\
        $S1$    & $231.0$   & $52.7$   & $1.14$   & $7.60$   & $5.70$   & $1.90$    & $0.75$           & $3.26$\\
        $S2$    & $216.0$   & $36.9$   & $1.08$   & $6.65$   & $5.38$   & $1.27$    & $0.81$           & $2.48$\\
        $S3$    & $209.7$   & $35.9$   & $1.05$   & $6.26$   & $5.20$   & $1.03$    & $0.85$           & $2.12$\\
        \bottomrule
    \end{tabular}
    
    \label{tab: friction values Reb 5300}
\end{table}
\begin{figure}
    \centering
    \begin{subfigure}[t]{0.495\textwidth}
        \centering
        \subcaption{}
        %\input{plots/Results/}
        \input{plots/Results/5300/5300Reb_comparison_uplus.pgf}
        \label{subfig: Reb 5300 velocity profile}
    \end{subfigure}
    \hfill
    \begin{subfigure}[t]{0.49\textwidth}
        \centering
        \subcaption{}
        \input{plots/Results/5300/5300Reb_velocityDefect.pgf}
        \label{subfig: Reb 5300 velocity defect}        
    \end{subfigure}
    \plotverticalspacing
    \begin{subfigure}[t]{0.495\textwidth}
        \centering
        \subcaption{}
        \input{plots/Results/5300/5300Reb_comparison_totalStress.pgf}
        \label{subfig: Reb 5300 total stresses}
    \end{subfigure}
    \hfill
    \begin{subfigure}[t]{0.495\textwidth}
        \centering
        \subcaption{}
        \input{plots/Results/5300/5300Reb_comparison_RS.pgf}
        \label{subfig: Reb 5300 reynolds stresses}
    \end{subfigure}
    \\[-10pt]
    \caption[Comparison of $S1$, $S2$ and $S3$ at $Re_b = 5300$]{Comparison between the tested surfaces at $Re_b = 5300$. \sOneFiveThreeSymbol \ denotes surface S1, \sTwoFiveThreeSymbol \ surface S2, \sThreeFiveThreeSymbol \ surface S3 and \referenceSymbol \ the reference data from~\cite{yaoDirectNumericalSimulations2023}.~(a) Mean velocity profiles.~(b) Velocity defect of local velocity to centreline velocity.~(c) \lsTotalStressSymbol \ total stress,\lsViscousStressSymbol \ Viscous stress, \lsReynoldsStressSymbol \ turbulent shear stress and \lsDispersiveStressSymbol \ dispersive stress.~(d) Reynolds stresses. \lsRuuSymbol \ $u'u'$, \lsRvvSymbol \ $v'v'$, \lsRwwSymbol \ $w'w'$ and \lsRuvSymbol \ $u'v'$}
    \label{fig: comp Reb 5300}
\end{figure}

Various mean friction and drag quantities for the rough surfaces are listed in \autoref{tab: friction values Reb 5300}.
Furthermore, \autoref{fig: comp Reb 5300} presents first- and second-order statistics for all three surfaces. 
In \autoref{subfig: Reb 5300 velocity profile}, the mean velocity profiles $U^+$ are displayed over the wall-normal distance $y^+$. 
Evidently, all three surfaces exhibit analogous behaviour within the log-law region of the flow, which further supports Townsend's similarity hypothesis.
This is in line with the findings of the velocity defects in \autoref{subfig: Reb 5300 velocity defect}, since all three curves collapse onto the reference data in the outer region.
The surface $S3$ attains the maximum $U^+$ value, $S1$ attains the minimum, while $S2$ is situated in between.
Due to the largest roughness height of surface $S1$, the flow experiences the most drag, hence the lowest mean velocity $U^+$.
The increased drag of $S1$ is attributable to both increased viscous drag $\tau_v$ and increased pressure drag $\tau_p$, in comparison with the other surfaces.
However, the influence of viscous drag on total drag is found to decrease for higher roughness height $k^+$ and effective slope $ES$, as shown in \autoref{tab: friction values Reb 5300}. 
Since $S3$ has a very similar surface topography but lower effective slope $ES$ compared to $S2$, the higher maximum mean velocity $U^+$ and thus reduced drag is not surprising.
Furthermore, the analysis indicates that $S3$ demonstrates the most significant influence of viscous drag on total drag, or conversely, the least significant influence of pressure drag.
In the near-wall region, the velocity behaviour varies significantly between the surfaces.
For all three surfaces the viscous sublayer is only weakly pronounced in the plots.
This is due to the fact that all the lines show the characteristic curve in the near-wall region, but with varying curvature.
Yet, they approach a positive $U^+$ value for $y^+ = 0$.
This outcome is, however, expected.
The $y^+ = 0$ position is at the constant radius $R = 1$.
Roughness valleys will exhibit radii greater than $1$.
%The $y^+ = 0$ position is within the flow for large parts of the pipe, depending on the respective wall location.
For all positions at which the wall is further away from the pipe centre than $y^+ = 0$, hence at negative $y^+$ regions, fluid flow is present at the virtual origin $y^+ = 0$ and beneath.
This results in positive $U^+$ values at the virtual origin. 
This effect is more pronounced for surface $S1$, since it exhibits the largest roughness height $k^+$, hence more room for fluid flow below $y^+ = 0$.

The total stress over the wall-normal position in outer units $r/R$ is plotted in \autoref{subfig: Reb 5300 total stresses}.
Its individual components, i.e. turbulent shear stress, viscous stress and dispersive stress are displayed as well.
It is evident that the expected slope of $1$ is matched for all surfaces in the entire domain, except in the near wall region. 
The total stress begins to drop at approximately $r/R = 0.15$, even though $S1$ declines slightly earlier.
The drop of the total stress is due to the viscous stress not reaching the anticipated value of $1$ at the virtual origin, as it would be expected for a smooth pipe.
The reasoning for that is analogous to that provided in the analysis of the velocity profile in the near-wall region.
This effect is most pronounced for $S1$, while it is weaker for $S2$ and $S3$.
The viscous stress behaves very similarly for $S2$ and $S3$, which indicates that the roughness height has the strongest impact on the viscous stress in rough conditions, while the effective slope only has a minor influence.
In the outer region, the viscous shear stress collapses on the smooth-wall curve and reaches zero at the pipe centre.
The turbulent stress displays analogous behaviour.
In the outer region, there is excellent agreement with the smooth-wall reference.
Moreover, the overall trend towards the wall is found to be highly comparable between the rough-wall cases and the smooth-wall case.
However, the rough-wall cases show smaller maximum turbulent stresses, with $S1$ reaching the lowest value.
At the wall, none of the three rough cases reach zero turbulent stress, remaining at values between $0.05$ and $0.15$.
Higher roughness heights result in increased turbulent stresses at the virtual origin, for the same reasons as for the velocity at the virtual origin.
The decreased turbulent stress in the rough cases is balanced by the dispersive stress.
The data indicates a positive correlation between the dispersive stresses and roughness height, with an additional influence of effective slope.
$S1$ shows the highest values of the dispersive stresses, with the peak sitting at $r/R \approx 0.1$.
The trend for $S2$ is the same but shifted slightly to lower values.
The maximum dispersive stresses achieved by $S3$ are comparable to those of $S2$, but with an outward shift in the maximum value to $r/R \approx 0.15$.
\autoref{subfig: Reb 5300 reynolds stresses} depicts the Reynolds stresses over the wall-normal distance in inner units.
The streamwise components $u'u'$ shows a significantly decreased peak values for all rough surfaces.
The data indicates a correlation between the $u'u'$ downshift and roughness height. 
The influence of the effective slope appears to vary with the wall-normal position. 
In the near-wall region, $S2$ and $S3$ deviate significantly, while the $u'u'$ value is nearly the same around the peak at approximately $y^+ = 20$.
Furthermore, for all rough surfaces the streamwise Reynolds stress peak is shifted to higher $y^+$ values when compared to the smooth wall pipe.
This shift occurs because the near-wall cycle, which is responsible for the streamwise Reynolds stress peak, is propagated outwards by the roughness~\citep{macdonaldTurbulentFlowTransitionally2016}.
The azimuthal Reynolds stresses $v'v'$ are increased over the entire domain, in comparison to the smooth cases.
This is the case for all three surfaces.
However, as in the streamwise direction, the maximum azimuthal Reynolds stresses are attained at higher $y^+$ values.
The same findings can be applied to the wall-normal Reynolds stresses $w'w'$.
The surface roughness exhibits only a minor influence on the Reynolds shear stress $u'v'$ with slightly higher values in the outer region.
The trends observed for the Reynolds stresses are due the disruption of turbulence by the roughness features.
This leads to the streamwise fluctuations being decreased.
Furthermore, a redistribution of turbulent kinetic energy away from the streamwise component is present, while increasing the azimuthal and wall-normal components.
This leads to the increased values for $v'v'$ and $w'w'$ ZITAT.

\begin{figure}
    \centering
    \begin{subfigure}{0.495\textwidth}
    \subcaption{}
        \subimport{plots/Results/5300/}{5300_spectra_piro.pgf}
    \end{subfigure}
    \hfill
    \begin{subfigure}{0.495\textwidth}
        \subcaption{}
        \subimport{plots/Results/5300/}{5300_spectra_kDown.pgf}
    \end{subfigure}
    \centering
    \begin{subfigure}{0.495\textwidth}
        \subcaption{}
        \subimport{plots/Results/5300/}{5300_spectra_scaled.pgf}
    \end{subfigure}
    \centering
    \\[5pt]
    \begin{subfigure}{\textwidth}
        \centering
        \hspace{0.9cm}
        \subimport{plots/Results/5300/}{3DSpectra_colorbar.pgf}
    \end{subfigure}
    \caption[Streamwise energy spectra at $Re_b = 5300$]{Streamwise pre-multiplied Energy spectra at $Re_b = 5300$. Wall-normal position plotted over wavelength. Colour denoting energy content, with blue indicating low energy content, and yellow high energy content. Red marker indicates the position of maximal energy. Contour lines start at 0.25 with intervals of 0.5.~(a) $S1$.~(b) $S2$.~(c) $S3$}
    \label{fig: Reb5300 energy spectra}
\end{figure}

For all three surfaces, the pre-multiplied energy spectra are plotted in \autoref{fig: Reb5300 energy spectra}.
It should be noted that the figure displays solely the non-dimensionalised energy spectra for the streamwise turbulent fluctuations, hence $k_zE_{uu}^+$.
The wall-normal position $y^+$ is plotted over the streamwise wavelength $\lambda_z^+$ in viscous units.
The energy content is denoted by the colour, as well as by contour lines.
Blue and purple shadings denote weak energy content, while green and yellow indicate
 strong energy content.
The contour lines start at $0.25$, with steps of $0.5$.
% The vertical first contourline in all three plots at $\lambda_z^+ \approx 200$ is due to the mesh resolution and the smallest wavelenghts that can be resolved by the grid point spacing.
A close examination of the positions of the energy peaks reveal a clear correlation between the peak energy content and surface roughness. 
The data indicates that an increase in roughness results in a shift of the energy peak outwards, away from the wall.
This peak value for $S1$ is at $y^+ \approx 23.0$, for $S2$ at $y^+ \approx 19.7$ and for $S3$ at $y^+ \approx 19.0$.
This is in line with the findings from the Reynolds stresses in \autoref{subfig: 5300validation Reynolds stress}.
Furthermore, the wavelength of the energy peak varies between the surfaces. 
For $S1$ the peak is located at $\lambda_z^+ \approx 528$ and for $S3$ at $\approx 752$, while it is located at $\lambda_z^+ \approx 542$ for $S2$.
This phenomenon may be attributed to the fact that the surfaces used for $S1$ and $S3$ exhibit larger wavelengths, or conversely, that $S2$ is generated using shorter wavelengths in order to be able to achieve the same roughness height with an unmodified effective slope.
Given this fact, the $S2$ energy peak located at a lower $\lambda_z^+$ value is in line with the findings of~\cite{chanSecondaryMotionTurbulent2018}.
Their data also suggest, that a lower wavelength of the surface roughness, leads to the energy peak being located at lower wavelengths.
%This is confirmed with the data of the current study, since the energy peak of $S2$ is located at a lower wavelength, than of $S3$.
% Their data also show a correlation between surface roughness wavelength and wavelength position of the energy peak, however for regular surface roughness.
% Additionally, their finding of the energy concentration at the roughness wavelenghts within the roughness canopy seems to be reinforced by the data of the current study.
The surface $S2$ shows a near-wall energy concentration at regions where $y^+ \approx 1$ at shorter wavelengths than for surfaces $S1$ and $S3$.
Also, the overall near-wall energy content is lower for $S2$, then for $S1$ and $S3$.
Both of these findings correspond to the findings by~\cite{chanSecondaryMotionTurbulent2018}, which indicates similar behaviour of the energy spectra for regular and irregular rough surfaces.
Additionally, the data of the current study suggests, that a smaller roughness height leads to the energy being distributed over a larger wavelength range.
However, this could also be due to the overall higher energy content.
% However, there is one main contradiction to the data of~\citeauthor{chanSecondaryMotionTurbulent2018}
% The highest energy peak can be observed for $S1$, hence the surface with the largest roughness height $k^+$ at $4.27$.
% $S2$ exhibits by far the lowest energy at $2.95$ and $S3$ sits in between at $3.97$.
% This is the opposite behaviour encountered by~\citeauthor{chanSecondaryMotionTurbulent2018}.
% The reason for the varying results remains to be investigated. 
% \begin{figure}
%     \centering
%     \begin{subfigure}[t]{0.495\textwidth}
%         \centering
%         \subcaption{}
%         \subimport{plots/Results/5300/}{5300_piro_instantaneousField_y+15.pgf}
%     \end{subfigure}
%     \hfill
%     \begin{subfigure}[t]{0.495\textwidth}
%         \centering
%         \subcaption{}
%         \subimport{plots/Results/5300/}{5300_piro_instantaneousField_y+115.pgf}
%     \end{subfigure}
%     \begin{subfigure}[t]{0.495\textwidth}
%         \centering
%         \subcaption{}
%         \subimport{plots/Results/5300/}{5300_kDown_instantaneousField_y+15.pgf}
%     \end{subfigure}
%     \hfill
%     \begin{subfigure}[t]{0.495\textwidth}
%         \centering
%         \subcaption{}
%         \subimport{plots/Results/5300/}{5300_kDown_instantaneousField_y+108.pgf}
%     \end{subfigure}
%     \begin{subfigure}[t]{0.495\textwidth}
%         \centering
%         \subcaption{}
%         \subimport{plots/Results/5300/}{5300_scaled_instantaneousField_y+15.pgf}
%     \end{subfigure}
%     \begin{subfigure}[t]{0.495\textwidth}
%         \centering
%         \subcaption{}
%         \subimport{plots/Results/5300/}{5300_scaled_instantaneousField_y+104.pgf}
%     \end{subfigure}
%     \begin{subfigure}{0.495\textwidth}
%         \centering
%         \subimport{plots/Results/5300/}{5300Reb_instantaneousField_colorbar_nearwall.pgf}
%     \end{subfigure}
%     \begin{subfigure}{0.495\textwidth}
%         \centering
%         \subimport{plots/Results/5300/}{5300Reb_instantaneousField_colorbar_centre.pgf}
%     \end{subfigure}
%     \caption{Streamwise velocity fluctuations of instantaneous flow fields. (a) and (b) Surface 1. (c) and (d) Surface 2. (e) and (f) Surface 3. (a), (c) and (e) at $y+ = 15$. (b), (d) and (f) at $y/R = 0.5$}
%     \label{fig: Reb 5300 instantaneous fields}
% \end{figure}





% Max values of Energy spectra:
% \begin{itemize}
%     \item 5300
%     \begin{itemize}
%         \item piro: 4.42
%         \item scaled: 3.97
%         \item kDown: 2.95
%     \end{itemize}
% \end{itemize}

\begin{figure}
    \centering
    \begin{subfigure}{0.75\textwidth}
        \centering
        \subimport{plots/Results/5300/}{5300Reb_pirotauWall_over_height.pgf}
    \end{subfigure}
    \caption[Trends of physical roughness height and viscous drag over streamwise coordinate]{Trends of physical roughness height and viscous drag over streamwise coordinate. Depicted for an exemplary $\theta$ position of the pipe.}
    \label{fig: roughness height and viscous stress Reb 5300 piro}
\end{figure}
\begin{figure}
    \centering
    \begin{subfigure}{0.7\textwidth}
        \hspace{0.9cm}
        \subimport{plots/Results/5300/}{5300Reb_piro_tauWall_over_height_bySlope.pgf}
    \end{subfigure}
    \caption[Scatter plot of viscous drag over physical roughness height]{Scatter plot of viscous drag over physical roughness height. Colour depicting the effective slope at the respective point}
    \label{fig: scatterplot piro Reb 5300}
\end{figure}

In the following, the connection between the viscous drag $\tau_v$ and the roughness shape is to be investigated by means of \autoref{fig: roughness height and viscous stress Reb 5300 piro} and \autoref{fig: scatterplot piro Reb 5300}.
It should be noted that the focus of this analysis is exclusively on surface $S1$, since the trends and conclusions can be directly applied to surfaces $S2$ and $S3$.
The plots with corresponding data of $S2$ and $S3$ are provided in \autoref{apx: viscous drag to surface height}.
The development of the viscous drag $\tau_v$ and the wall position $h$ over the z-coordinate at an arbitrarily chosen $\theta$ position of the pipe is shown in \autoref{fig: roughness height and viscous stress Reb 5300 piro}.
%The height position is normalised with the mean radius, hence the height deviation is shown.
For both the roughness height and the viscous drag, the fluctuations around their respective mean are plotted.
The viscous drag is non-dimensionalised into wall units.
\autoref{fig: roughness height and viscous stress Reb 5300 piro} indicates a clear correlation between  surface height and viscous drag.
This verifies the findings by~\cite{maDirectNumericalSimulation2021}.
The correlation appears to be stronger for positive slopes.
However, this is merely one sample of the surface.
To show the same influence across the entire surface, a scatter plot is constructed in \autoref{fig: scatterplot piro Reb 5300}.
Each data point on the scatter plot corresponds to a single interpolation point.
The number of interpolation points equals the number of grid points of the simulation mesh, but they are distributed uniformly in space.
The scatter plot illustrates the viscous drag in wall units over the surface height, employing the same conventions used in \autoref{fig: roughness height and viscous stress Reb 5300 piro}.
The colour map indicates the streamwise slope at the respective point.
Red shadings denote negative streamwise slopes while blue shadings denote positive streamwise slopes.
It is evident that a strong correlation exists between the viscous drag and the surface height for positive streamwise slopes, due to the straight band of blue points that is consistently rising from negative surface heights up to the highest points of the surface.
This shows that the correlation found in \autoref{fig: roughness height and viscous stress Reb 5300 piro} is not a local anomaly but rather, it is present for the entire surface.
However, a pronounced cluster of points is evident at a negative viscous drag over the entire range of surface heights.
The majority of these points are coloured deep red, hence indicating a strong negative streamwise slope.
A potential reason for this behaviour may be attributed to flow separation and recirculation due to the steep gradient which is oriented away from the flow.
Furthermore, the data indicates that negative slopes result in reduced overall local viscous drag, since most red points are in the lower half of the $\tau_v$ values, while the upper half of the plot is mostly occupied by blue data points.
In addition, there are instances where both a positive slope and negative viscous drag are observed. 
All of these points are situated in proximity to the mean radius, or in some cases beneath it.
These points may demonstrate the same circulations in roughness valleys as the negative sloped points, albeit at the other end where the surface is rising again.

\fxinfo{ab hier neu}
\section{Simulations with 11700$Re_b$}\label{sec: results 11700 Reb}

In this section, the simulation results for the Reynolds numbers $Re_b = 11700$ are presented.
As done in \autoref{sec: results 5300Reb}, a comparison was made between the three surfaces.
The smooth reference data is taken from~\cite{elkhouryDirectNumericalSimulation2013}.
\begin{table}
    \centering
    \caption[Computed friction values for $Re_b = 11700$]{Computed friction values for $Re_b = 11700$}
    \begin{tabular}{c c c c c c c c c}
        \toprule
                & $Re_\tau$ & $k^+$     & $C_f\ [\times10^{-3}]$ & $\tau_w\ [\times10^{-3}]$    & $\tau_v\ [\times10^{-3}]$ & $\tau_p\ [\times10^{-3}]$ & $\tau_v/\tau_w$   & $\Delta U^+$\\
        \midrule
        smooth  & $360.1$   & -         & $7.58$                & $3.79$                        & $3.79$                    & -                         & $1.00$            & - \\
        $S1$    & $487.4$   & $111.1$   & $9.19$                & $6.94$                        & $4.60$                    & $2.35$                    & $0.66$            & $4.05$\\
        $S2$    & $453.9$   & $77.6$    & $8.81$                & $6.02$                        & $4.40$                    & $1.62$                    & $0.73$            & $3.02$\\
        $S3$    & $431.9$   & $73.9$    & $8.44$                & $5.45$                        & $4.22$                    & $1.23$                    & $0.77$            & $3.81$\\
        \bottomrule
    \end{tabular}
    \label{tab: friction values 11700 Reb}
\end{table}
\begin{figure}
    \centering
    \begin{subfigure}[t]{0.495\textwidth}
        \centering
        \subcaption{}
        %\input{plots/Results/}
        \input{plots/Results/11700/11700Reb_comparison_uplus.pgf}
        \label{subfig: Reb 11700 velocity profile}
    \end{subfigure}
    \hfill
    \begin{subfigure}[t]{0.49\textwidth}
        \centering
        \subcaption{}
        \input{plots/Results/11700/11700Reb_velocityDefect.pgf}
        \label{subfig: Reb 11700 velocity defect}        
    \end{subfigure}
    \plotverticalspacing
    \begin{subfigure}[t]{0.495\textwidth}
        \centering
        \subcaption{}
        \input{plots/Results/11700/11700Reb_comparison_totalStress.pgf}
        \label{subfig: Reb 11700 total stresses}
    \end{subfigure}
    \hfill
    \begin{subfigure}[t]{0.495\textwidth}
        \centering
        \subcaption{}
        \input{plots/Results/11700/11700Reb_comparison_RS.pgf}
        \label{subfig: Reb 11700 reynolds stresses}
    \end{subfigure}
    \\[-10pt]
    \caption[Comparison of $S1$, $S2$ and $S3$ at $Re_b = 5300$]{Comparison between the  tested surfaces at $Re_b = 5300$. \sOneElevenSevenSymbol \ denotes surface S1, \sTwoElevenSevenSymbol \ surface S2, \sThreeElevenSevenSymbol \ surface S3 and \referenceSymbol \ the reference data from~\cite{elkhouryDirectNumericalSimulation2013}.~(a) Mean velocity profiles.~(b) Velocity defect of local velocity to centreline velocity.~(c) \lsTotalStressSymbol \ total stress,\lsViscousStressSymbol \ Viscous stress, \lsReynoldsStressSymbol \ turbulent shear stress and \lsDispersiveStressSymbol \ dispersive stress.~(d) Reynolds stresses. \lsRuuSymbol \ $u'u'$, \lsRvvSymbol \ $v'v'$, \lsRwwSymbol \ $w'w'$ and \lsRuvSymbol \ $u'v'$}
    \label{fig: comp Reb 11700}
\end{figure}

Various mean friction and drag quantities for the rough surfaces at $Re_b = 11700$ are listed in 
\autoref{tab: friction values 11700 Reb}.
First- and second-order statistics are presented in \autoref{fig: comp Reb 11700}.
In \autoref{subfig: Reb 11700 velocity profile}, the mean velocity profiles $U^+$ over the wall-normal distance $y^+$ are plotted.
The velocity profiles demonstrate a comparable trend to those observed in cases with  $Re_b = 5300$.
Increased roughness height leads to an increase in the\fxinfo{Aussage etwas anders: prüfen,} roughness functions.
Furthermore, $S2$ and $S3$ show almost the same velocity at the virtual origin, while the velocity of $S2$ is significantly lower.
The velocity defects are shwon in \autoref{subfig: Reb 11700 velocity defect}.  They display  excellent agreement in the proximity of the pipe centre.
Overall, $S3$ closely mirrors a smooth function\fxinfo{Satz war unvollständig, prüfen}, with the exception of the region near the wall.
$S1$ and $S2$ behave very similarly throughout the domain, exhibiting an overall elevated trend from\fxinfo{when?} $r/R > 0.05$.



The stress budgets are depicted in \autoref{subfig: Reb 11700 total stresses}.
It is evident that the total stress displays the anticipated behaviour within  the entire domain, as it aligns with the expected slope of $1$ for all three surfaces.
The total stress begins to decline at the earliest stage for $S1$, which  is consistent with the behaviour observed for $Re_b = 5300$.
The turbulent stresses undergo a collapse  in the outer region when confronted with the smooth wall case.
Surfaces $S2$ and $S3$ show nearly identical behaviour across the entire domain.
In the outer region, there is concurrence  with the smooth wall case, yet they reach their peak at lower values.
The turbulent stress exhibited by $S1$ deviates significantly.
As expected, the data shows the correct slope near the centre. However,   from $r/R \approx 0.75$ onwards, the turbulent stresses are subtantially reduced.
Additionally, a decisive shift in the slope of the curve is evident.
This behaviour is compensated by the dispersive stress of $S1$, which shows a nearly constant trend from the wall until $r/R \approx 0.4$ and reaches zero at $r/R \approx 0.7$.
The Reynolds stresses in \autoref{subfig: Reb 11700 reynolds stresses} show overall  trends as they  have been expected.

$S1$ exhibits the lowest streamwise Reynolds stress peak, conversely, the Reynolds stresses are stronger for $S2$ and $S3$.
Furthermore, the outward shift of the Reynolds stress peak is more pronounced for larger roughness heights in this data set.
The azimuthal Reynolds stresses $v'v'$ behave almost identically between the three surfaces, except in the near wall region.In this region,  larger roughness height result in larger  azimuthal fluctuations.
The same holds for the wall-normal Reynolds stresses $w'w'$.
The Reynolds shear stresses $u'v'$ exhibit higher values in the outer region when compared to the smooth reference data.
Nevertheless, a decline in the range is evident, from $y^+ \approx 10$ to $y^+ \approx 100$.
Surfaces $S2$ and $S3$ attain and surpass the $u'v'$ value of the smooth case at an earlier point,  at $y^+ \approx 80$. This corresponds approximately to the  values of the roughness height in wall units $k^+$.
This is due to the dispersive stresses reducing the turbulent stresses within the roughness  canopy.

\begin{figure}
    \centering
    \begin{subfigure}{0.495\textwidth}
        \subcaption{}
        \subimport{plots/Results/11700/}{11700_spectra_piro.pgf}
        \label{subfig: 11700 spectra piro}
    \end{subfigure}
    \hfill
        \begin{subfigure}{0.495\textwidth}
            \subcaption{}
            \subimport{plots/Results/11700/}{11700_spectra_kDown.pgf}
            \label{subfig: 11700 spectra kDown}
    \end{subfigure}
    \\[-20pt]
    \begin{subfigure}{0.495\textwidth}
        \subcaption{}
        \subimport{plots/Results/11700/}{11700_spectra_scaled.pgf}
        \label{subfig: 11700 spectra scaled}
    \end{subfigure}
    \centering
    \begin{subfigure}{\textwidth}
        \centering
        \hspace{0.9cm}
        \subimport{plots/Results/11700/}{3DSpectra_colorbar.pgf}
    \end{subfigure}
    \caption[Streamwise pre-multiplied energy spectra at $Re_b = 11700$]{Streamwise pre-multiplied Energy spectra at $Re_b = 11700$. Wall-normal position plotted over wavelength. Colour denoting energy content, with blue indicating low energy content, and yellow high energy content. Red marker indicates the position of maximal energy. Contour lines start at 0.25 with intervals of 0.5.~(a) $S1$.~(b) $S2$.~(c) $S3$}
    \label{fig: energy spectra 11700}
\end{figure}


The plots of the energy spectra of the simulations using $Re_b = 11700$ are provided in \autoref{fig: energy spectra 11700}.
The maximum energy content for $S1$ is $1.00$, for $S2$ it is $1.21$, and for $S3$ it is $1.27$.
As the roughness height increases, the energy peak observed to be shifted outwards. This phenomenon is evident  at $y^+ \approx 35.2$ for $S1$, at $27.1$ for $S2$ and at $y^+ \approx 25.6$ for $S3$.
Like for the $Re_b = 5300$ cases, this corresponds to  the position of the $u'u'$ peak.
Additionally, the dominant wavelength is found to be longer for smaller roughness heights, since it is $\lambda_z^+ \approx 557$ for $S1$ and $\lambda_z^+ \approx 638$.
The dominant wavelength  for $S2$ is $\lambda_z^+ \approx 543$.
This shortening for $S2$ compared to the similar roughness height at surface $S3$ can be explained by the surface geometry of $S2$. The surface of $S2$  is generated using of a greater number of Fourier modes compared to $S3$.
Due to the smaller roughness elements, the turbulence is divided in smaller  sections,  consequently leading to shorter energy wavelengths.
Furthermore, the energy is distributed over a wider wavelength range  for $S2$ and $S3$ in comparison to for $S1$.
In particular, regions at high wavelengths and high $y^+$ values are more pronounced.
Additionally, it is evident that the contour line at $0.75$ encloses a larger area in the plots for $S2$ and $S3$ than for $S1$. This also implies a higher energy content overall.
The relatively vertical contour line at the higher wavelength limit of \autoref{subfig: 11700 spectra scaled} can be attributed to the maximum resolvable wavelength by the simulation domain.
With a $\delta_\nu = 0.002315$ and a pipe length of $L_z = 8\pi$, the maximum resolvable wavelength is $L_z / \delta_\nu \approx 10850$.
\begin{figure}
    \centering
    \begin{subfigure}{0.75\textwidth}
        \centering
        \subimport{plots/Results/11700/}{11700Reb_pirotauWall_over_height.pgf}
    \end{subfigure}
    \caption[Trends of physical roughness height and viscous drag over streamwise coordinate for $Re_b = 11700$]{Trends of physical roughness height and viscous drag over streamwise coordinate. Depicted for an exemplary $\theta$ position of the pipe of $S1$ at $Re_b = 11700$.}
    \label{fig: roughness height and viscous stress Reb 11700 piro}
\end{figure}
\begin{figure}
    \centering
    \begin{subfigure}{0.7\textwidth}
        \hspace{0.9cm}
        \subimport{plots/Results/11700/}{11700Reb_piro_tauWall_over_height_bySlope.pgf}
    \end{subfigure}
    \caption[Scatter plot of viscous drag over physical roughness height at $Re_b = 11700$]{Scatter plot of viscous drag over physical roughness height. Colour depicting the effective slope at the respective point. Data shown for $S1$ at $Re_b = 11700$.}
    \label{fig: scatterplot piro Reb 11700}
\end{figure}


The development of the viscous drag $\tau_v$ with respect to the streamwise position and the roughness height is plotted in \autoref{fig: roughness height and viscous stress Reb 11700 piro}.
The scatter plot in \autoref{fig: scatterplot piro Reb 11700} is analogous to the scatter plot in  \autoref{fig: scatterplot piro Reb 5300} in \autoref{sec: results 5300Reb}.
As done in \autoref{sec: results 5300Reb}, the plot for $S1$ solely is provided in this section. The data for $S2$ and $S3$  are  incorporated in \autoref{apx: viscous drag to surface height}.
When examining the development of the viscous drag and the roughness height of a single $\theta$ position, the positive correlation between the two quantities is confirmed for a $Re_b = 11700$.
Furthermore, the plot suggests that the viscous drag tends to decrease sharply near the roughness crest and, in many cases, slightly before the crest.
This phenomenon may be indicative of  flow detachment near roughness peaks, particularly within  the positive $h/\delta$ region.
\autoref{fig: scatterplot piro Reb 11700} depicts analogous relations as \autoref{fig: scatterplot piro Reb 5300}.
A positive correlation is demonstrated between the viscous drag and effective streamwise slope for positive slopes.
The finding regarding flow detachment in \autoref{fig: roughness height and viscous stress Reb 11700 piro} is confirmed by the scatter plot, which shows that this behaviour is present over the whole surface.
However, the $Re_b = 11700$ case exhibits  a tendency towards lower viscous drag values, which consistent with the  anticipation for growing Reynolds numbers.
This effect is also reflected by the lower friction coefficients $C_f$ compared to the $Re_b = 5300$ cases.

\section{Simulations with 19000$Re_b$}\label{sec: results 19000 Reb}

This section is devoted to the presentation of the results for the simulations using $Re_b = 19000$.

\begin{table}
    \centering
    \caption[Computed friction values for $Re_b = 19000$]{Computed friction values for $Re_b = 19000$}
    \begin{tabular}{c c c c c c c c c}
        \toprule
                & $Re_\tau$ & $k^+$     & $C_f\ [\times10^{-3}]$    & $\tau_w\ [\times10^{-3}]$    & $\tau_v\ [\times10^{-3}]$ & $\tau_p\ [\times10^{-3}]$ & $\tau_v/\tau_w$   & $\Delta U^+$\\
        \midrule
        smooth  & $547.8$   & -         & $6.65$                    & $3.32$                        & $3.32$                    & -                         & $1.00$            & - \\
        $S1$    & $763.0$   & $174.0$   & $7.95$                    & $6.45$                        & $3.98$                    & $2.47$                    & $0.62$            & $5.61$\\
        $S2$    & $715.7$   & $122.4$   & $7.71$                    & $5.68$                        & $3.86$                    & $1.82$                    & $0.68$            & $3.98$\\
        $S3$    & $670.8$   & $114.7$   & $7.37$                    & $4.99$                        & $3.68$                    & $1.30$                    & $0.73$            & $3.74$\\
        \bottomrule
    \end{tabular}
    \label{tab: friction values 19000 Reb}
\end{table}


\autoref{tab: friction values 19000 Reb} shows the most significant values associated with the friction behaviour of the surfaces.
Similar behavioural patterns are exhibited for $Re_b = 19000$ compared to $Re_b = 11700$ and $Re_b = 5300$.
The skin friction coefficient $C_f$ declines with a decrese in both  roughness height and  effective slope. But it is significantly increased in rough conditions.
Both the viscous drag and the pressure drag decrease for decreasing roughness height and effective slope.
Furthermore, the influence of viscous drag on total drag increases for decreasing roughness height.
The friction function decreases with decreasing roughness height and effective slope, a finding that is consistent with the other results presented in this table.
\begin{figure}
    \centering
    \begin{subfigure}[t]{0.495\textwidth}
        \centering
        \subcaption{}
        %\input{plots/Results/}
        \input{plots/Results/19000/19000Reb_comparison_uplus.pgf}
        \label{subfig: Reb 19000 velocity profile}
    \end{subfigure}
    \hfill
    \begin{subfigure}[t]{0.49\textwidth}
        \centering
        \subcaption{}
        \input{plots/Results/19000/19000Reb_velocityDefect.pgf}
        \label{subfig: Reb 19000 velocity defect}        
    \end{subfigure}
    \plotverticalspacing
    \begin{subfigure}[t]{0.495\textwidth}
        \centering
        \subcaption{}
        \input{plots/Results/19000/19000Reb_comparison_totalStress.pgf}
        \label{subfig: Reb 19000 total stresses}
    \end{subfigure}
    \hfill
    \begin{subfigure}[t]{0.495\textwidth}
        \centering
        \subcaption{}
        \input{plots/Results/19000/19000Reb_comparison_RS.pgf}
        \label{subfig: Reb 19000 reynolds stresses}
    \end{subfigure}
    \\[-10pt]
    \caption[Comparison of $S1$, $S2$ and $S3$ at $Re_b = 19000$]{Comparison between the 3 tested surfaces at $Re_b = 19000$. \sOneNineteenSymbol \ denotes surface S1, \sTwoNineteenSymbol \ surface S2, \sThreeNineteenSymbol \ surface S3 and \referenceSymbol \ the reference data from~\cite{yaoDirectNumericalSimulations2023}.~(a) Mean velocity profiles.~(b) Velocity defect of local velocity to centreline velocity.~(c) \lsTotalStressSymbol \ total stress,\lsViscousStressSymbol \ Viscous stress, \lsReynoldsStressSymbol \ turbulent shear stress and \lsDispersiveStressSymbol \ dispersive stress.~(d) Reynolds stresses. \lsRuuSymbol \ $u'u'$, \lsRvvSymbol \ $v'v'$, \lsRwwSymbol \ $w'w'$ and \lsRuvSymbol \ $u'v'$}
    \label{fig: comp Reb 19000}
\end{figure}


The first- and second-order statistics are presented in \autoref{fig: comp Reb 19000}.
The mean velocity profiles are plotted in \autoref{subfig: 19000 validation velocity profile}.
The mean velocities for $S1$ and $S2$ collapse onto another at the virtual origin. This does not imply convergence for $S1$ and $S2$ in that regard.
Rather, it can be observed that the velocities at the virtual origin move closer together across the three Reynolds numbers considered in this thesis.
The velocity defects illustrated in \autoref{subfig: Reb 19000 velocity defect} are in accordance with in the outer region.
In the proximity of  the wall, $S1$ and $S2$ collapse onto another, while $S3$ exhibits  a decline in values.
The stress budget is plotted in \autoref{subfig: 19000 validation total stress}.
%The total stress budget are plotted in \autoref{subfig: 19000 validation total stress}.
The total stresses demonstrate a high degree of agreement  between the surfaces.
It is evident that all  three surfaces show the anticipated slope of $1$ throughout the domain, except in the near wall region.
However, turbulent stresses manifest the expected slope of $1$ solely in the outer region.
From a radial position of $r/R < 0.8$, the trends of the turbulent stresses are decreased noticeably compared to the smooth wall reference.
This effect is found to be more pronounced  for  the surface $S1$  with  greater roughness height.
The peak turbulent stresses are reduced significantly for all three surfaces when compared to the smooth wall reference.
%The dispersive stresses exhibit a distinct trend.
The dispersive stresses for $S1$ and $S3$ display  a similar trend, however, with a lower magnitude for $S3$.  This is attributable  to the lower roughness height but similar surface topography.
$S2$ exhibits an almost linear decrease of the dispersive stress until reaching zero at $r/R \approx 0.7$.
The viscous stresses demonstrate only a minor  influence at this Reynolds number.
The values remain close to  zero until $r/R \approx 0.1$ for all three surfaces, and rise to values between $0.21$ for $S1$ and $0.34$ for $S3$ at the virtual origin.
The findings for the Reynolds stresses are displayed in \autoref{subfig: 19000 validation Reynolds stress}. Reynolds stress values correspond to the $Re_b = 11700$ and $Re_b = 5300$ cases.
The streamwise Reynolds stresses $u'u'$ decrease for rising roughness heights and effective slopes. 
Furthermore, the peak is shifted outwards.
The overall increase of the spanwise $v'v'$ and wall-normal $w'w'$ fluctuations is also present.


\begin{figure}
    \centering
    \begin{subfigure}{0.495\textwidth}
        \subcaption{}
        \subimport{plots/Results/19000/}{19000_spectra_piro.pgf}
        \label{subfig: 11700 spectra piro}
    \end{subfigure}
    \hfill
        \begin{subfigure}{0.495\textwidth}
            \subcaption{}
            \subimport{plots/Results/19000/}{19000_spectra_kDown.pgf}
            \label{subfig: 11700 spectra kDown}
    \end{subfigure}
    \\[-20pt]
    \begin{subfigure}{0.495\textwidth}
        \subcaption{}
        \subimport{plots/Results/19000/}{19000_spectra_scaled.pgf}
        \label{subfig: 11700 spectra scaled}
    \end{subfigure}
    \centering
    \begin{subfigure}{\textwidth}
        \centering
        \hspace{0.9cm}
        \subimport{plots/Results/19000/}{3DSpectra_colorbar.pgf}
    \end{subfigure}
    \caption[Streamwise pre-multiplied energy spectra at $Re_b = 19000$]{Streamwise pre-multiplied Energy spectra at $Re_b = 19000$. Wall-normal position plotted over wavelength. Colour denoting energy content, with blue indicating low energy content, and yellow high energy content. Red marker indicates the position of maximal energy. Contour lines start at 0.25 with intervals of 0.5.~(a) $S1$.~(b) $S2$.~(c) $S3$}
    \label{fig: energy spectra 19000}
\end{figure}

The pre-multiplied streamwise energy spectra in \autoref{fig: energy spectra 19000} also show analogous trends to those observed for $Re_b = 11700$ and $Re_b = 5300$.
The maximum energy content decreases with roughness height and effective slope, since for $S1$  it is at $0.88$, For $S2$ at $1.04$ and for $S3$ at $1.11$.
In addition, the positions of the energy peaks demonstrate a similar behaviour to that observed for $Re_b = 11700$ and $Re_b = 5300$.
The energy peak for $S1$ is located at $y^+ \approx 45.5$, for $S2$ at $y^+ \approx 33.7$ and for $S3$ at $y^+ \approx 31.4$.
However, the behaviour-dominant wavelengths are distinct for $Re_b = 190000$.
$S1$ manifests the longest dominant wavelength at $\lambda_z^+ \approx 710$, while it is significantly shorter for $S2$ at $\lambda_z^+ \approx 600$ and also shortened for $S3$ at $\lambda_z^+ \approx 674$. 
INTERPRETIEREN

\begin{figure}
    \centering
    \begin{subfigure}{0.75\textwidth}
        \centering
        \subimport{plots/Results/19000/}{19000Reb_pirotauWall_over_height.pgf}
    \end{subfigure}
    \caption[Trends of physical roughness height and viscous drag over streamwise coordinate at $Re_b = 19000$]{Trends of physical roughness height and viscous drag over streamwise coordinate. Depicted for an exemplary $\theta$ position of the pipe. Data for $S1$ and $Re_b = 19000$}
    \label{fig: roughness height and viscous stress Reb 19000 piro}
\end{figure}
\begin{figure}
    \centering
    \begin{subfigure}{0.7\textwidth}
        \hspace{0.9cm}
        \subimport{plots/Results/19000/}{19000Reb_piro_tauWall_over_height_bySlope.pgf}
    \end{subfigure}
    \caption[Scatter plot of viscous drag over physical roughness height at $Re_b = 19000$]{Scatter plot of viscous drag over physical roughness height. Colour depicting the effective slope at the respective point. Data shown for $S1$ at $Re_b = 19000$}
    \label{fig: scatterplot piro Reb 19000}
\end{figure}


The development of the viscous drag with respect to the roughness height for an arbitrary postion of the parameter $\theta$ is plotted in \autoref{fig: roughness height and viscous stress Reb 19000 piro}.
Overall the findings from $Re_b = 11700$ and $Re_b = 5300$ are confirmed for the case of $Re_b = 19000$.
A positive slope leads to steeply rising viscous drag values.
Furthermore, the positive correlation between the roughness height and viscous drag is also validated.
For the plot of the entire surface in \autoref{fig: scatterplot piro Reb 19000}, the analogous trends  are observed.
The values for the viscous drag are found to be lower in comparison to  the values of $Re_b = 11700$ and $Re_b = 5300$.
Due to this extent, the conclusions remain unaltered.

\section{Comparison between Reynolds numbers}
In the following section, the influence of the Reynolds number is investigated.
To this extent, the data from \autoref{sec: results 5300Reb}, \autoref{sec: results 11700 Reb} and \autoref{sec: results 19000 Reb} are taken, and the same surfaces are compared against one another.
\subsection{Reynolds number influence on surface $S1$} \label{subsec: Reynolds comp S1}

In this subsection, the behaviour of the surface $S1$ is compared for the three different Reynolds numbers.
\begin{figure}
    \centering
    \begin{subfigure}{0.495\textwidth}
        \centering
        \subcaption{}
        \input{plots/Results/comp_Reynolds/piro/Reynolds_comparison_uplus_piro.pgf}
        \label{subfig: Reynolds comp S1 velocity}
    \end{subfigure}
    \hfill
    \begin{subfigure}{0.495\textwidth}
        \centering
        \subcaption{}
        \input{plots/Results/comp_Reynolds/piro/Reynolds_comparison_totalStress_piro.pgf}
        \label{subfig: Reynolds comp S1 total stress}
    \end{subfigure}
    \plotverticalspacing
    \begin{subfigure}{0.495\textwidth}
        \centering
        \subcaption{}
        \input{plots/Results/comp_Reynolds/piro/Reynolds_comparison_RS_piro_uu.pgf}
        \label{subfig: Reynolds comp S1 Reynolds stresses uu}
    \end{subfigure}
    \hfill
    \begin{subfigure}{0.495\textwidth}
        \centering
        \subcaption{}
        \input{plots/Results/comp_Reynolds/piro/Reynolds_comparison_RS_piro_vv.pgf}
        \label{subfig: Reynolds comp S1 Reynolds stresses vv}
    \end{subfigure} 
    \\[-10pt]  
    \caption[Comparison of surface $S1$ with different Reynolds numbers]{Comparison between Reynolds numbers for surface $S1$.  \sOneFiveThreeSymbol\ denotes $Re_b = 5300$, \sOneElevenSevenSymbol \ $Re_b = 11700$ and \sOneNineteenSymbol\ $Re_b = 19000$.~(a) Mean velocity profiles.~(b) \lsTotalStressSymbol \ total stress, \lsReynoldsStressSymbol \ turbulent shear stress and \lsDispersiveStressSymbol \ viscous stress.~(c) Reynolds stresses. \lsRuuSymbol \ $u'u'$ and \lsRuvSymbol \ $u'v'$.~(d) Reynolds stresses. \lsRvvSymbol \ $v'v'$ and \lsRwwSymbol \ $w'w'$}
    \label{fig: results Reynolds comparison S1}
\end{figure}


In \autoref{fig: results Reynolds comparison S1} the first- and second-order statistics are plotted.
To improve the clarity of the plots, some data lines are omitted.
In \autoref{subfig: Reynolds comp S1 total stress} the dispersive stresses are neglected, since the total stresses, indicated by the dashed lines, are the sum of the  viscous stresses, turbulent shear stresses and dispersive stresses.
Both viscous stresses and turbulent shear stresses are shown, indicating a sufficient amount of information.
The plot of the Reynolds stresses is split up into two plots.
The streamwise Reynolds stresses $u'u'$ and the Reynolds shear stresses $u'v'$ are depicted in \autoref{subfig: Reynolds comp S1 Reynolds stresses uu}.
The  azimuthal stresses $v'v'$ and wall-normal stresses $w'w'$ are plotted in \autoref{subfig: Reynolds comp S1 Reynolds stresses vv}.


The mean velocity profiles in \autoref{subfig: Reynolds comp S1 velocity} show a clear trend.
The velocity at the virtual origin $y^+ = 0$ is increasing with rising Reynolds number.
%This indicates that the streamwise component of the turbulent flow is able to penetrate deeper into the roughness indentations at higher Reynolds numbers.
Furthermore, an increasing downshift of the log-law, and thus an increasing roughness function, can be determined with an increasing Reynolds number.
It is noteworthy that that the maximum non-dimensionalised velocity $U^+$ is exhibits a comparable value  for all three Reynolds numbers. 
This finding aligns with the results reported of ~\cite{demaioDirectNumericalSimulation2023}.
The total stresses in \autoref{subfig: Reynolds comp S1 total stress} behave almost the same between the Reynolds numbers.
The maximum total stress increases with increasing Reynolds number, and the maximum value for the turbulent stress shifts  towards the wall with increasing Reynolds number.


The viscous stresses and turbulent stresses decrease monotonically with increasing Reynolds number, which in turn implies increasing dispersive stresses.
This effect becomes particularly pronounced for $Re_b = 19000$, which is not unexpected, given the particular high roughness height $k^+$ in this case.
Furthermore, the viscous stresses decrease with increasing Reynolds numbers.
This  may be attributed to two distinct phenomena.
First, the viscous drag decreases naturally with increasing Reynolds number.
Second, the reduced maximum viscous stress can be associated with the increased velocity at $y^+ = 0$, as shown in \autoref{subfig: Reynolds comp S1 velocity}.
\autoref{subfig: Reynolds comp S1 Reynolds stresses uu} depicts the Reynolds stresses for the three Reynolds numbers.
The streamwise Reynolds stresses $u'u'$ demonstrate a monotonic shift towards lower peak values at larger $y^+$ values.
Furthermore, the peak becomes less pronounced for increasing Reynolds numbers, since the stresses reach a value at $y^+ = 0$ that is only slightly lower than the peak value.
At $Re_b = 19000$, the outer peak of the streamwise fluctuations $u'u'$ starts to emerge.
The turbulent fluctuations $u'v'$ show similar behaviour.
In particular, $Re_b = 11700$ and $Re_b = 19000$ show a significant overlap between $y^+ = 40$ and $y^+ = 100$.
The azimuthal fluctuations $v'v'$ and the wall-normal fluctuations $w'w'$ exhibit a high degree of similarity in their behaviour.
With an increasing Reynolds number, the maximum fluctuations also increase.
Furthermore, the location of the peak is trending towards higher $y^+$ values.
The reasoning for this phenomenon is very similar to that given in \autoref{sec: results 5300Reb}.
%However, the data demonstrates, that this effect strengthens with increasing Reynolds number.


\subsection{Reynolds number influence on surface $S2$}\label{subsec: Reynolds comp S2}

\begin{figure}
    \centering
    \begin{subfigure}{0.495\textwidth}
        \centering
        \subcaption{}
        \input{plots/Results/comp_Reynolds/kDown/Reynolds_comparison_uplus_kDown.pgf}
        \label{subfig: Reynolds comp S2 velocity}
    \end{subfigure}
    \hfill
    \begin{subfigure}{0.495\textwidth}
        \centering
        \subcaption{}
        \input{plots/Results/comp_Reynolds/kDown/Reynolds_comparison_totalStress_kDown.pgf}
        \label{subfig: Reynolds comp S2 total stress}
    \end{subfigure}
    \plotverticalspacing
    \begin{subfigure}{0.495\textwidth}
        \centering
        \subcaption{}
        \input{plots/Results/comp_Reynolds/kDown/Reynolds_comparison_RS_kDown_uu.pgf}
        \label{subfig: Reynolds comp S2 Reynolds stresses uu}
    \end{subfigure}
    \hfill
    \begin{subfigure}{0.495\textwidth}
        \centering
        \subcaption{}
        \input{plots/Results/comp_Reynolds/kDown/Reynolds_comparison_RS_kDown_vv.pgf}
        \label{subfig: Reynolds comp S2 Reynolds stresses vv}
    \end{subfigure} 
    \\[-10pt]  
    \caption[Comparison of surface $S2$ with different Reynolds numbers]{Comparison between Reynolds numbers for surface $S2$.  \sTwoFiveThreeSymbol\ denotes $Re_b = 5300$, \sTwoElevenSevenSymbol \ $Re_b = 11700$ and \sTwoNineteenSymbol\ $Re_b = 19000$.~(a) Mean velocity profiles.~(b) \lsTotalStressSymbol \ total stress, \lsReynoldsStressSymbol \ turbulent shear stress and \lsDispersiveStressSymbol \ viscous stress.~(c) Reynolds stresses. \lsRuuSymbol \ $u'u'$ and \lsRuvSymbol \ $u'v'$.~(d) Reynolds stresses. \lsRvvSymbol \ $v'v'$ and \lsRwwSymbol \ $w'w'$}
    \label{fig: results Reynolds comparison S2}
\end{figure}


The first- and second-order statistics of surface $S2$ at the three tested Reynolds numbers are compared in \autoref{fig: results Reynolds comparison S2}.
The mean velocity profiles in \autoref{subfig: Reynolds comp S2 velocity} show the clear, yet anticipated  trend of an increased downshift of the log-law for an increasing Reynolds number, thereby implying an increased roughness function.
As observed for the case of $S1$, the maximum velocity $U^+$ is approximately equivalent  for all three Reynolds numbers.
Furthermore, the velocity at the virtual origin increases with increasing Reynolds number.
The stress budgets are plotted in \autoref{subfig: Reynolds comp S2 total stress}.
The total stress peak is at higher values for an increasing Reynolds number.
It is also shifted towards the wall.
Contrary to $S1$, the turbulent stress for $Re_b = 5300$ and $Re_b = 11700$ exhibits equivalent behaviour  until $r/R \approx 0.4$.
Beyond that point, the turbulent stress of $Re_b = 5300$ starts to decline while it keeps growing with the constant slope until $r/R \approx 0.2$.
At $Re_b = 19000$ the turbulent stress behaves quite differently.
The graph shows a lower slope from $r/R \approx 0.6$ until $r/R \approx 0.27$, at which point the slope increases again. This results in the turbulent stress reaching a slightly lower maximum than in the $Re_b = 11700$ case at $r/R \approx 0.1$.
This indicates a substantially increased dispersive stress influence for $Re_b = 19000$.

Furthermore, the data suggests that the influence of dispersive stress does not increase linearly with Reynolds number.
In \autoref{subfig: Reynolds comp S3 Reynolds stresses uu} the streamwise Reynolds stresses and shear stresses are plotted.
The peak of the streamwise fluctuations $u'u'$ decreases for an increasing Reynolds number, accompanied by a shift outwards.
However, the difference between the peak values increases, compared to the higher roughness surface $S1$.
The shear stresses $u'v'$ show similar peak values for all three Reynolds numbers, exhibiting a consistent overall behaviour.
The spanwise and wall-normal fluctuations $v'v'$ and $w'w'$, respectively,  exhibit the same trends across the Reynolds numbers for $S2$ as they do for $S1$, as shown in \autoref{subfig: Reynolds comp S3 Reynolds stresses vv} 
The peaks increase with rising Reynolds numbers, while  simultaneously shifting towards higher values of  $y^+$.



\subsection{Reynolds number influence on surface $S3$}\label{subsec: Reynolds comp S3}

The results for the simulations using surface $S3$ at the three Reynolds numbers are compared.
\begin{figure}
    \centering
    \begin{subfigure}{0.495\textwidth}
        \centering
        \subcaption{}
        \input{plots/Results/comp_Reynolds/scaled/Reynolds_comparison_uplus_scaled.pgf}
        \label{subfig: Reynolds comp S3 velocity}
    \end{subfigure}
    \hfill
    \begin{subfigure}{0.495\textwidth}
        \centering
        \subcaption{}
        \input{plots/Results/comp_Reynolds/scaled/Reynolds_comparison_totalStress_scaled.pgf}
        \label{subfig: Reynolds comp S3 total stress}
    \end{subfigure}
    \plotverticalspacing
    \begin{subfigure}{0.495\textwidth}
        \centering
        \subcaption{}
        \input{plots/Results/comp_Reynolds/scaled/Reynolds_comparison_RS_scaled_uu.pgf}
        \label{subfig: Reynolds comp S3 Reynolds stresses uu}
    \end{subfigure}
    \hfill
    \begin{subfigure}{0.495\textwidth}
        \centering
        \subcaption{}
        \input{plots/Results/comp_Reynolds/scaled/Reynolds_comparison_RS_scaled_vv.pgf}
        \label{subfig: Reynolds comp S3 Reynolds stresses vv}
    \end{subfigure} 
    \\[-10pt]  
    \caption[Comparison of surface $S3$ with different Reynolds numbers]{Comparison between Reynolds numbers for surface $S3$.  \sThreeFiveThreeSymbol\ denotes $Re_b = 5300$, \sThreeElevenSevenSymbol \ $Re_b = 11700$ and \sThreeNineteenSymbol\ $Re_b = 19000$.~(a) Mean velocity profiles.~(b) \lsTotalStressSymbol \ total stress, \lsReynoldsStressSymbol \ turbulent shear stress and \lsDispersiveStressSymbol \ viscous stress.~(c) Reynolds stresses. \lsRuuSymbol \ $u'u'$ and \lsRuvSymbol \ $u'v'$.~(d) Reynolds stresses. \lsRvvSymbol \ $v'v'$ and \lsRwwSymbol \ $w'w'$}
    \label{fig: results Reynolds comparison S3}
\end{figure}


In \autoref{fig: results Reynolds comparison S3} the first- and second-order statistics are presented.
The velocity profiles in \autoref{subfig: Reynolds comp S2 velocity} manifest  a comparable behaviour to that exhibited by surfaces $S1$ and $S2$.
The increased roughness function becomes apparent due to the downshift of the log-law with increasing strength with increasing Reynolds number.
However, the maximum non-dimensionalised velocity $U^+$ increases with rising Reynolds number, thereby suggesting a correlation with the lower roughness for $S3$ and this particular observation\fxinfo{ mir unklar}.
The stress budgets are presented in \autoref{subfig: Reynolds comp S3 total stress}.
As for the other surfaces, the peak of the total stress rises with increasing Reynolds number, which implicates the peak being located closer to the wall, since the slope for all cases is $1$.
The turbulent stresses show comparable behaviour for $Re_b = 5300$ and $Re_b = 11700$.
However, at $Re_b = 11700$, a higher maximum turbulent stress is reached.
At $Re_b = 19000$, a pronounced impact of dispersive stress is apparent since the turbulent stress is significantly decreased from $r/R \approx 0.75$ until it reaches approximately the same peak value as $Re_b = 11700$. However, it is\fxinfo{prüfen} closer towards the wall.

This finding signifies an erratic escalation in  the influence of the dispersive stresses between $Re_b = 11700$ and $Re_b = 19000$.
The viscous stresses for $Re_b = 11700$ and $Re_b = 19000$ show almost identical behaviour from the pipe centre until $r/R \approx 0.3$.
Beyond that point, the viscous influence at $Re_b = 11700$ is higher than at $Re_b = 19000$.
It is evident that the viscous stresses at $Re_b = 5300$ are stronger throughout the domain.
The streamwise Reynolds stresses $u'u'$ for $S3$  are presented in \autoref{subfig: Reynolds comp S3 Reynolds stresses uu}. They demonstrate analogous trends to those observed for $S1$ and $S2$.
The peak experiences a downshift for increasing Reynolds numbers, as well as being shifted towards higher $y^+$ values.
However, this effect is more pronounced as for $S1$ and\fxinfo{??} for $S2$. THis further supports the argument that larger roughness heights lead to a stronger Reynolds number influence on streamwise fluctuations.
Furthermore, the streamwise fluctuations increase at the virtual origin, as they do at $S1$ and $S2$.
The shear stresses $u'v'$ behave analogously to $S2$.
As do the spanwise fluctuations $v'v'$ and wall-normal fluctuations $w'w'$.

% \section{Comparison between \acrshort{ibm} and body-fitted mesh}

% In this section a brief comparison between similar cases with different mesh applications is provided.

% \begin{figure}
%     \centering
%     \input{plots/Validation/4400_validation_Uplus.pgf}
%     \caption[Comparison of Velocity profiles of \acrshort{ibm} and body-fitted mesh]{Comparison between the mean velocity profiles of an \acrshort{ibm} and body-fitted case. Both simulations performed using $Re_b = 4400$ and similar roughness topographies. Data for \acrshort{ibm} taken from~\cite{demaioDirectNumericalSimulation2023}}
%     \label{fig: ibm comp}
% \end{figure}
% In \autoref{fig: ibm comp} the mean velocity profiles of a simulation performed using an \acrshort{ibm} mesh and one using a body-fitted mesh is plotted.
% The data of the \acrshort{ibm} mesh is taken from~\cite{demaioDirectNumericalSimulation2023}.
% For the body-fitted mesh is the same as used for the $Re_b = 5300$ cases of the current study.
% Both simulations are performed using $Re_b = 4400$.
% The surface is the grit blasted of~\cite{demaioDirectNumericalSimulation2023} and surface $S1$ of the current study.
% \begin{table}
%     \centering
%     \begin{tabular}{c c c c c c c}
%                         & $S_a$     & $S_q$     & $Sk$      $S_{z_{max}}$       & $Es$      & $S_{z_{5x5}}$\\
%         \acrshort{ibm}  & $0.031$   & \\
%         $S1$            & $0.031$
%     \end{tabular}
% \end{table}
% The surface quantities are
